\chapter{Análisis Numérico y Fundamentos del FEM}
\label{app:calculo}

Este apéndice recoge los conceptos de análisis numérico y teoría de la aproximación que sustentan la confiabilidad de las simulaciones presentadas en el libro. Su objetivo es proporcionar al lector ingeniero las herramientas necesarias para interpretar errores, estimar convergencia y comprender las decisiones de discretización tomadas en cada capítulo.

%==================================================================================
\section{Análisis del Error}
%==================================================================================

La validez de cualquier resultado numérico depende de cuantificar la distancia entre la solución exacta $u$ y su aproximación computacional $u_h$. Se definen dos medidas fundamentales:

\begin{definicion}[Error absoluto y relativo]
    Sea $x \in \R$ un valor exacto y $\tilde{x} \in \R$ su aproximación numérica. El \textbf{error absoluto} se define como:
    \begin{equation}
        e_{\text{abs}} = |x - \tilde{x}|.
    \end{equation}
    Si $x \neq 0$, el \textbf{error relativo} se define como:
    \begin{equation}
        e_{\text{rel}} = \frac{|x - \tilde{x}|}{|x|}.
    \end{equation}
\end{definicion}

En el contexto del Método de Elementos Finitos, el error se mide en normas funcionales sobre los espacios de Sobolev. Para una función $v$ definida sobre $\Omega$, las normas vectoriales de referencia son:
\begin{align}
    \norm{v}_{\ell_1} &= \sum_{i=1}^{n} |v_i|, \\
    \norm{v}_{\ell_2} &= \left( \sum_{i=1}^{n} |v_i|^2 \right)^{1/2}, \\
    \norm{v}_{\ell_\infty} &= \max_{1 \leq i \leq n} |v_i|.
\end{align}

Para matrices $A \in \R^{n \times n}$, la norma inducida asociada a $\norm{\cdot}_{\ell_p}$ se define como $\norm{A}_p = \sup_{v \neq 0} \norm{Av}_p / \norm{v}_p$.

\begin{notatecnica}
En el problema de la tubería del Capítulo~1, la elección de \texttt{mesh\_size=0.005} no es arbitraria. Para elementos Lagrange P2, la teoría del FEM garantiza convergencia de orden $\mathcal{O}(h^3)$ en la norma $L^2$ para problemas elípticos con solución suficientemente regular. Reducir $h$ por un factor de 2 disminuye el error en $L^2$ aproximadamente por un factor de 8, aunque el coste computacional crece como $\mathcal{O}(h^{-2})$ en 2D.
\end{notatecnica}

%==================================================================================
\section{Condición del Problema}
%==================================================================================

La sensibilidad de la solución de un sistema lineal $A\bm{x} = \bm{b}$ ante perturbaciones en los datos viene cuantificada por el \textbf{número de condición} de la matriz.

\begin{definicion}[Número de condición]
    Sea $A \in \R^{n \times n}$ una matriz invertible. El número de condición respecto a una norma matricial $\norm{\cdot}$ se define como:
    \begin{equation}
        \kappa(A) = \norm{A} \, \norm{A^{-1}}.
    \end{equation}
    En particular, para la norma espectral $\norm{\cdot}_2$, se tiene $\kappa_2(A) = \sigma_{\max} / \sigma_{\min}$, donde $\sigma$ denota los valores singulares de $A$.
\end{definicion}

Si $\kappa(A) \gg 1$, pequeñas perturbaciones en $\bm{b}$ o en $A$ pueden inducir cambios arbitrariamente grandes en $\bm{x}$. En tales casos, el problema se dice \textbf{mal condicionado} (ill-conditioned).

En el caso industrial de la tubería de concentrado, el contraste de conductividades térmicas entre el acero y el concentrado es:
\begin{equation}
    \frac{k_{\text{acero}}}{k_{\text{conc}}} = \frac{45.0}{2.0} = 22.5.
\end{equation}
Este salto de coeficientes introduce heterogeneidad en la matriz de rigidez, pero en un problema bidimensional con mallas moderadas no genera un mal condicionamiento severo que impida la resolución directa. No obstante, en geometrías 3D industriales con millones de grados de libertad y contrastes de material de dos órdenes de magnitud, el uso de un \textbf{precondicionador} (por ejemplo, AMG o ILU) se vuelve imprescindible para mantener la eficiencia del solver iterativo.

%==================================================================================
\section{Cuadratura Gaussiana}
%==================================================================================

En la implementación computacional del FEM, las integrales de la forma débil no se evalúan analíticamente, sino que se aproximan mediante \textbf{cuadratura numérica}. En FEniCSx, el operador \texttt{ufl.dx} desencadena automáticamente una regla de cuadratura adaptada al espacio de elementos finitos seleccionado.

Para elementos Lagrange P2 en 2D, los términos de rigidez involucran productos de gradientes de funciones cuadráticas, lo que da lugar a integrandos polinomiales de grado hasta 4 en cada dirección. La \textbf{cuadratura de Gauss-Legendre} con 3 puntos por dirección integra exactamente polinomios de grado hasta 5, lo cual resulta \textbf{suficiente} para garantizar que el error de cuadratura no domine sobre el error de discretización del método.

\begin{notatecnica}
El usuario avanzado puede sobrescribir la regla de cuadratura en FEniCSx mediante \texttt{ufl.dx(metadata=\{\ldots\})} si requiere integración exacta de términos no polinomiales (por ejemplo, funciones de material definidas tabularmente) o si desea reducir el coste computacional mediante cuadratura reducida en problemas no lineales.
\end{notatecnica}
